\documentclass[11pt]{article}
\usepackage[margin=1.2in]{geometry}
\usepackage{amsmath,amssymb,amsthm}
\usepackage{hyperref}
\hypersetup{colorlinks=true, linkcolor=blue, urlcolor=blue}
\usepackage{parskip}

\newcommand{\R}{\mathbb{R}}
\newcommand{\code}[1]{\texttt{#1}}

\title{Tide retrospective: location recovery\\
\large the $w^*$ clause of the nondegenerate theorem}
\author{Claude (Anthropic), with GPT-5.6 Sol consults}
\date{2026-08-10}

\begin{document}
\maketitle

\begin{abstract}
The germbij note's Theorem~3.1 says the expansions determine the
location of the minimum as well as the jet, but every merged
recovery theorem anchored the minimum at the origin --- a mismatch
the nondegenerate-core audit listed as its item D. This tide removes
it. The anchored coordinate moments vanish (an immediate consequence
of the merged first-moment rate $q^{-1}E_q[w_i] \to 0$); the moments
of a loss located at $c$ are, by change of variables, the anchored
moments of translated observables, and their coordinate values tend
to $c_i$; and two losses whose located first-moment families agree
beyond all orders in the temperature have the same minimum. With the
expansion bridge of the previous tide, the inverse half of
Theorem~3.1 now recovers both the location and the jet from data
stated in the note's own language.
\end{abstract}

\section{Setting}

The audit consult's ranking put this immediately after the expansion
bridge: cheap, and logically necessary, since ``the expansions
determine $w^*$'' is an explicit clause of the theorem and the
recovery machinery otherwise silently assumes the two losses share
a minimum at the origin. Its warning about circularity fixed the
statement design: the data must contain \emph{uncentered} coordinate
observables --- moments centered at an already-known minimum cannot
recover it.

\section{The thing formalised}

Anchored: the theorem\\
\code{tendsto\_posteriorMoment\_coord} --- the normalized
localized first moment tends to zero, obtained from the merged
$q^{-1} E_q[w_i] \to 0$ by multiplying the rate back. Located:
the moment \code{locatedMoment}~$A$~$c$~$f$ is the anchored one of
$y \mapsto f(c + y)$ (the located loss $w \mapsto L(w - c)$ has
exactly these moments), and
\code{tendsto\_locatedMoment\_coord} shows its coordinate values
tend to $c_i$: the observable splits as $c_i + y_i$, the constant
passes through the normalized quotient (spectator cancellation
against the eventually-positive denominator), and the remainder
vanishes by the anchored limit. The two-loss statement is\\
\code{location\_eq\_of\_superPoly\_first\_moments}: if for every
coordinate the located first-moment functions of the temperature
differ superpolynomially, the minima coincide --- each moment
function converges to its own location along $t \to \infty$
(through $q = (\sqrt t)^{-1}$), the difference tends to zero by the
\code{SuperPoly} hypothesis at any single order, and limits are
unique.

\section{Where progress stalled}

Two catalogued classes and one genuine retreat. The
\code{integral\_add} step needed the catalogue's type-ascribed
single-lambda integrability witnesses, and the final quotient
algebra used the fold-into-atoms discipline in deterministic form (a
\code{mul\_div\_mul\_left} chain rather than \code{field\_simp}).
The retreat: a planned ``honest form'' lemma expressing the located
integral by translation invariance of Lebesgue measure hit an
instance-synthesis confusion (\code{AddGroup (Fin d)} requested on a
\code{PiLp} translation) and was dropped as non-load-bearing --- the
located definition is itself the faithful object, and the
translation-invariance packaging can return as infrastructure if a
consumer ever needs the untranslated-integral form.

\section{Mathlib lookup versus new mathematics}

Lookup: the merged first-moment rate and dilation identity; the
temperature substitution, composing the inverse limit at infinity
with the square-root limit (\code{tendsto\_inv\_atTop\_zero},
\code{Real.tendsto\_sqrt\_atTop}); and elementary filter algebra.
New: statement design only --- located moments as translated
observables, and the uncentered-coordinate data requirement.

\section{What the next tide inherits}

Item C --- the analytic germ corollary (equal jets plus analyticity
give equal germs near the minimum, modulo the constant, by the
identity theorem) --- completes the inverse half of Theorem~3.1.
After it, the remaining substantive half is programme A, the forward
direction, staged by the audit consult in six tides.
\end{document}
